\chapter{方法と確認状態}
CSVを正本とし、索引・頻度表・問題ブロックを派生出力する。本文の手書き部分と自動生成部分は保存先を分ける。
生OCRは未校正の抽出値として保持し、校正済みの条件・選択肢対応・図参照を別レコードに置く。
確認証跡には原典ハッシュと対象内容のハッシュを記録する。原典が変われば旧確認を引き継がない。
\section{統計の限界}
暫定分類による数と、確認済みの部分集合の数を並べる。確認済みの部分集合は全出題頻度の推定値ではない。
難易度・学習時間が空欄なら未評価であり、0や易問として処理しない。
\section{完成の条件}
実際の校正・解説、他肢の検証、引用・参照、TeXビルド、専門査読が必要である。
ラベルの存在とコンパイルの成功だけでは、教材の内容品質や完成を認定しない。
